% Generated from analysis.json by tools/write_values.py; CC-BY-4.0, Jonathan Yantis.
\def\HaveCtwo{1}
\def\HaveWarm{1}
\newcommand{\AbstractResult}{The longest fully replayed rolling device frontier is 64 seconds on C-1, with 128 prefixes, P99 760 ns and maximum 879 ns in that window.}
\newcommand{\AbstractDuration}{A separately selected fixed mature 128-run cohort resolves 75 s passing (P99 983 ns) and 76 s failing on a one-second grid. Later capture ends vary: C-2 later ends at 24 s; C-1 later ends at 24 s.}
\newcommand{\PrimaryTable}{C-1 initial & $\geq48$ s & 690/827/911 ns & 128; 0 / 0\\
C-2 initial & $\geq48$ s & 504/606/880 ns & 128; 0 / 0\\
C-2 later & $24$ s & 727/942/1020 ns & 128; 1 / 0\\
C-1 later & $24$ s & 652/925/1267 ns & 128; 1 / 0\\}
\newcommand{\TableScope}{All rows use installed v33 and the latest per-age ring at the declared capture end. A passing tested frontier with no adequately populated next-age failure is a lower bound. These values do not describe every status in the window.}
\newcommand{\PrimaryResults}{Four fixed capture leases, each bounded at 600 s and stopped near 596 s, preserve separate device/time windows. The table reports declared capture-end rings. Across all captured 48 s prefixes in the initial windows, C-1 has 220 prefixes (P99/max 911/953 ns) and C-2 has 193 (796/880 ns), with no failures. C-2's later device briefly reports 64 s before the new segment's long-age prehistory is flushed; the lighter plot line retains that observation without promoting it to an independently reconstructed 64 s numeric ring. The conditioned 120 s follow-up was not triggered because final 96 s qualification failed. The fully replayed 64 s result spans 15 observed statuses at 2785.510--2925.768 s device uptime in C-1's later window.}
\newcommand{\CoverageResults}{The captures contain 25639 accepted pairs and 10450 numerically replayed START/END pairs, with zero missing accepted-pair sequence records or capture-loss counter deltas. Gap counts above 250 ms are 212/338/292/213; their maxima are 1.688/1.375/2.188/1.000 s, respectively in table order. Orphan, unfinished and unknown outcomes remain separate in the supplement.}
\newcommand{\RampResults}{C-2 acquisition is captured before trial 1; its first 48 s qualification is at 400.690 s uptime. C-1's initial setup refusal occurred before USB opening; the retry begins at trial 668 and independently flushes each age's history. Its exact device trajectory reaches 48 s, drops to 1.5 s at 480.719 s because short runs expire unknown, and recovers at 520.782 s. C-2's later window also has an unknown-driven short-age dip, followed by actual longer-age timing failures and recovery. These are distinct causes of a lower displayed frontier.}
\newcommand{\FineGridResults}{C-2 later's same fixed cohort gives 32 s on the coarse grid and 44 s on the one-second grid; 39 younger long starts remain censored separately. C-1 later's same fixed cohort gives 64 s on the coarse grid and 75 s on the one-second grid; 39 younger long starts remain censored separately.}
\newcommand{\FineGridTable}{C-2 later & 128 & 44 s & 45 s & 975 / 1023 & 2 / 0\\
C-1 later & 128 & 75 s & 76 s & 983 / 1028 & 2 / 0\\}
\newcommand{\FineGridScope}{P99 pairs in this table are at the last passing and first failing one-second ages, with the failure/unknown count at the failing age. All numeric prefixes at those boundaries are explicitly counted in the supplement. The same-cohort coarse comparison separates grid resolution from changing run populations; it does not substitute this older fixed cohort for the capture-end OLED ring.}
